\subsection{R-16: simultaneous chords give an unbounded genuine-padding range}

\paragraph{Statement R-16.}
For $m\geq4$ put
\[
 h_m=\min\{\lfloor m/2\rfloor,m-3\},\qquad
 F_m=h_m(m-3).
\]
Thus $h_4=1$ and $h_m=\lfloor m/2\rfloor$ for $m\geq5$.
The quaternary permanent image satisfies the uniform bound
\begin{equation}\label{R16:dimension}
 \dim\mathcal P_{m,4}\geq 2m^2+2+F_m.
\end{equation}
At $m=4$ and $m=5$, this lower bound is respectively $35$ and
$56$, the whole ambient dimension. Thus it also proves the
quartic and quintic dominance of R-12 by an algebraic argument.
Consequently, if $n>m\geq4$ and
\begin{equation}\label{R16:paddingrange}
 2n^2\leq2m^2+F_m+2,
\end{equation}
then, in every sufficiently large coordinate-ring degree $d$,
some partition $\lambda\vdash nd$ with exactly four rows has
\[
 m_\lambda\bigl(\ell^{n-m}\operatorname{per}_m\bigr)
   >m_\lambda(\det_n).
\]
In particular, $n=m+1$ satisfies this conclusion for every
$m\geq12$. Together with the finite certificates in R-14,
this proves the conclusion for $n=m+1$ and every $m\geq8$.
For every fixed real $1<\rho<\sqrt5/2$, the same conclusion
holds for $n=\lfloor\rho m\rfloor$ and all sufficiently large
$m$. These are multiplicity gaps in a specified range of
linear padding; the proof does not specify the partition or
give a polynomial-time procedure to find it.

\begin{proof}
Fix a primitive $m$th root of unity $\zeta$ and use indices
$0,\ldots,m-1$ modulo $m$. Specialize the cyclic pencil of
R-15 to
\[
 \ell_i=x-\zeta^i y,\qquad v_i=z-\zeta^i w.
\]
Let $M_0$ have diagonal $\ell_i$ and cycle entries
$(M_0)_{i,i+1}=v_i$. Its differential image is denoted $W_0$.
The restriction-line and interval-product proof of R-15 applies
to this specialization, since both lists of roots are distinct,
and gives $\dim W_0\geq2m^2+2$.

We will simultaneously deform the disjoint row pairs
\[
 \mathcal A=\{0,2,\ldots,2(h_m-1)\},\qquad
 \{a,a+1\}\quad(a\in\mathcal A).
\]
For each $a\in\mathcal A$ add the two entries
\begin{equation}\label{R16:simultaneouspencil}
 (M_\varepsilon)_{a,a+2}=p_a=\varepsilon\ell_a,
 \qquad
 (M_\varepsilon)_{a+1,a}=q_a=\varepsilon v_{a+2}.
\end{equation}
All other entries remain those of $M_0$. The pairs of rows are
disjoint, including when the last triple of column indices
wraps around the cycle. Consequently rows $a,a+1$ have support
only in the three columns $a,a+1,a+2$, even when every other
pair has also been perturbed.

For $r\notin\{a,a+1\}$ let $R_{a,r}(\varepsilon)$ be the
permanent obtained by deleting rows $a,a+1,r$ and columns
$a,a+1,a+2$. Write $C_{r,j}(\varepsilon)$ for the complementary
permanent deleting row $r$ and column $j$. In each of the
three cases $j=a,a+1,a+2$, the two rows $a,a+1$ must match
the other two columns of that triple: they have no entries
outside it, and exactly two of its columns remain. The rest
of the matching is arbitrary in the same complementary
$(m-3)\times(m-3)$ matrix. Therefore, with no hypothesis on
the other row pairs,
\[
 \begin{aligned}
 C_{r,a}&=(v_a v_{a+1}+p_a\ell_{a+1})R_{a,r},\\
 C_{r,a+1}&=(\ell_a v_{a+1}+p_aq_a)R_{a,r},\\
 C_{r,a+2}&=(\ell_a\ell_{a+1}+v_aq_a)R_{a,r}.
 \end{aligned}
\]
Here and in the next display all permanents are evaluated at
$\varepsilon$. Subtraction gives the exact identity
\begin{equation}\label{R16:simultaneousidentity}
 \ell_a C_{r,a}-v_a C_{r,a+1}-p_a C_{r,a+2}
   =-2\varepsilon^2\ell_a v_a v_{a+2}R_{a,r}(\varepsilon).
\end{equation}
It follows that the polynomial vector
\[
 T_{a,k}(\varepsilon)
   =\ell_a v_a v_{a+2}R_{a,a+k}(\varepsilon)
\]
belongs to the full differential image $W_\varepsilon$ for
every $\varepsilon\ne0$ and $2\leq k\leq m-2$.
At $\varepsilon=0$ the remaining cycle paths have a unique
matching, so
\begin{equation}\label{R16:normalvectors}
 T_{a,k}:=T_{a,k}(0)
 =\ell_a v_a v_{a+2}
   \left(\prod_{i=a+2}^{a+k-1}v_i\right)
   \left(\prod_{i\notin[a,a+k]}\ell_i\right).
\end{equation}
Intervals and products are cyclic; the first product is empty
when $k=2$. This polynomial has bidegree $(m-k,k)$ in
$(x,y)$ and $(z,w)$.

We next prove that all $h_m(m-3)$ vectors in
\eqref{R16:normalvectors} are linearly independent modulo
$W_0$. This is the step that permits simultaneous rank growth,
rather than adding ranks from different deformations.

For $1\leq k<m$ set
\[
 F_{a,k}
   =\left(\prod_{i=a}^{a+k-2}v_i\right)
     \left(\prod_{i\notin[a,a+k-1]}\ell_i\right).
\]
The bidegree-$(m-k,k)$ part of $W_0$ is exactly
\begin{equation}\label{R16:bidegreespace}
 (W_0)_{m-k,k}
   =\operatorname{span}\{zF_{a,k},wF_{a,k}:0\leq a<m\}.
\end{equation}
Indeed the cycle cofactors with $k-1$ factors $v_i$ give
these generators directly. Every other generator of the
same bidegree is a cofactor with $k$ such factors multiplied
by a linear form $u(x,y)$. Write
$u=A\ell_a+B\ell_{a+k}$, possible because these two
linear forms are independent. The resulting polynomial is
$A v_aF_{a+1,k}+B v_{a+k-1}F_{a,k}$, proving the reverse
inclusion. This accounts for every cofactor and multiplier
of that bidegree. The two endpoint bidegrees have dimensions
at most $m+1$, so this argument also gives
$\dim W_0\leq2(m+1)+2m(m-1)=2m^2+2$.

Let $\tau$ act by
\[
 \tau(x,y,z,w)=(x,\zeta y,z,\zeta w).
\]
Then $F_{a,k}=\tau^aF_{0,k}$ and
$T_{a,k}=\tau^aT_{0,k}$. Equation
\eqref{R16:bidegreespace} shows that $W_0$ is invariant
under this cyclic action. Let $\pi_t$ denote its character
projection with eigenvalue $\zeta^t$.
We claim that
\begin{equation}\label{R16:allmodes}
 \pi_t(T_{0,k})\notin\pi_t(W_0)
 \quad\text{for every }2\leq k,t\leq m-2.
\end{equation}

Here is the full coefficient calculation proving the claim.
Write
\[
 \begin{aligned}
 A(x,y)&=\prod_{i=k}^{m-1}(x-\zeta^iy)
       =\sum_r A_r x^{m-k-r}y^r,\\
 B(z,w)&=\prod_{i=0}^{k-2}(z-\zeta^iw),\\
 E(z,w)&=\prod_{i=1}^{k-2}(z-\zeta^iw),\qquad B=(z-w)E.
 \end{aligned}
\]
Use subscripts for coefficients in the powers of the second
variable, with coefficients outside the degree range defined
to be zero. Every $A_r$, $0\leq r\leq m-k$, is nonzero.
For completeness the finite geometric-product identity gives
\[
 A_r=(-1)^r\zeta^{kr+r(r-1)/2}
      \prod_{j=1}^{r}
       \frac{1-\zeta^{m-k-r+j}}{1-\zeta^j}.
\]
This identity follows by induction on the number of linear
factors. Every exponent in numerator or denominator lies
strictly between $0$ and $m$, so no factor vanishes.

We have $F_{0,k}=AB$. The $(x,y)$ factor of $T_{0,k}$ is
$A(x,\zeta y)$, since its roots are $k+1,\ldots,m-1,0$.
If its $(z,w)$ factor is $B'(z,w)$, then
\[
 \begin{aligned}
 D(z,w):=B'(z,\zeta^{-1}w)
   &=(z-\zeta^{-1}w)(z-\zeta w)E(z,w)\\
   &=(z-w)B(z,w)+c\,zwE(z,w),\\
 c&=2-\zeta-\zeta^{-1}\ne0.
 \end{aligned}
\]
These equalities hold also for $k=2$, when $E=1$.

In character $t$, with $2\leq t\leq m-2$, the monomials
are precisely those with $r+s=t$; a second possibility
$r+s=t+m$ exceeds total degree $m$.
The allowed indices $s$ form the interval
\[
 \max\{0,t-(m-k)\}\leq s\leq\min\{k,t\}.
\]
It contains at least three consecutive integers, because
$2\leq k,t\leq m-2$.
The coefficient of such a monomial in $T_{0,k}$ is
$\zeta^t A_{t-s}D_s$. The two possible generators for
the character-$t$ part of \eqref{R16:bidegreespace}, namely
the projections of $zF_{0,k}$ and $wF_{0,k}$, have
coefficients $A_{t-s}B_s$ and $A_{t-s}B_{s-1}$.
Since all the displayed $A_{t-s}$ are nonzero, membership
in their span would imply that $D_s$ is a fixed linear
combination of $B_s$ and $B_{s-1}$ on the allowed interval.
Subtracting $(z-w)B$ and using $c\ne0$, it would follow
that $E_{s-1}$ is also such a linear combination.

That last assertion is impossible on any three consecutive
allowed indices $s=j,j+1,j+2$. To see this, use
$B_s=E_s-E_{s-1}$ to perform invertible column operations:
the relevant three-column determinant is nonzero if and
only if
\begin{equation}\label{R16:toeplitz}
 \det\bigl(E_{j+u-v}\bigr)_{0\leq u,v\leq2}\ne0.
\end{equation}
Here $0\leq j\leq k-2$. Put $N=k-2$.
Since $E_i=(-1)^i e_i(\zeta,\ldots,\zeta^N)$, row and
column signs change \eqref{R16:toeplitz} to the determinant
\[
 \det\bigl(e_{j+u-v}(\zeta,\ldots,\zeta^N)\bigr)_{0\leq u,v\leq2}
   =s_{(3^j)}(\zeta,\ldots,\zeta^N),
\]
by the dual Jacobi--Trudi identity (its representation-ring
form is equation~(2.4) in the attached proof of Sawin's Lemma~2.3;
R-B15). This Schur value is nonzero. If $N=0$, necessarily $j=0$ and the determinant
is one. Otherwise use the Weyl alternant at
$(1,\zeta,\ldots,\zeta^{N-1})$, then multiply by the
nonzero scalar $\zeta^{3j}$. For
$\lambda=(3^j,0^{N-j})$, the exponents
$\lambda_i+N-i$ are distinct integers in $[0,N+2]=[0,k]$.
Since $k<m$, their powers of $\zeta$ are distinct. The
alternant numerator is therefore a nonzero Vandermonde
determinant. Its denominator has the distinct exponents
$N-i$ and is nonzero as well. This proves
\eqref{R16:toeplitz}, and hence \eqref{R16:allmodes}.

Now fix $k$ and suppose
$\sum_{a\in\mathcal A}u_aT_{a,k}\in W_0$.
Projecting to each of the characters in
\eqref{R16:allmodes} gives
\[
 \sum_{a\in\mathcal A}u_a\zeta^{at}=0
 \qquad(2\leq t\leq m-2).
\]
There are $h_m\leq m-3$ distinct indices $a$.
Use the $h_m$ consecutive equations
$t=2,\ldots,h_m+1$. Their matrix is an ordinary
Vandermonde matrix in the distinct numbers $\zeta^a$,
with each column multiplied by the nonzero factor
$\zeta^{2a}$. It is invertible, so every $u_a=0$.
Thus the $h_m$ vectors for this $k$ are independent
modulo $W_0$. Different $k$ have different bidegrees
and $W_0$ is bigraded, so all $h_m(m-3)$ vectors are
independent modulo $W_0$.

Choose $2m^2+2$ differential columns at $M_0$ forming a
basis of $W_0$, and continue the same columns at
$M_\varepsilon$. Adjoin all the polynomial vectors
$T_{a,k}(\varepsilon)$. At zero these columns are
independent by the preceding argument, so one coefficient
minor of size $2m^2+2+F_m$ is a polynomial in
$\varepsilon$ with nonzero constant term. For generic
nonzero $\varepsilon$ it remains nonzero, and all its
adjoined vectors lie in $W_\varepsilon$ by
\eqref{R16:simultaneousidentity}. Hence
$\dim W_\varepsilon\geq2m^2+2+F_m$.
Differential rank bounds image dimension from below,
proving \eqref{R16:dimension}.

For padding, R-14 gives
$\dim\mathcal Z_{m,n}=\dim\mathcal P_{m,4}+3$.
Under \eqref{R16:paddingrange} this dimension is at least
$2n^2+3$, strictly larger than the determinant upper
bound $2n^2+2$. The Hilbert-function and highest-weight
argument of R-14 therefore supplies the asserted
exactly-four-row multiplicity gaps in every sufficiently
large degree. The explicit degree threshold of G-10
also applies to these pairs.

For $m\geq12$ one has
$F_m\geq(m-1)(m-3)/2\geq4m$, so
\eqref{R16:paddingrange} holds at $n=m+1$.
R-14 already covers each remaining $8\leq m\leq11$.
Finally, $F_m=m^2/2+O(m)$, and the right side of
\eqref{R16:paddingrange}, divided by $2m^2$, tends to
$5/4$. This proves the stated asymptotic range for every
fixed $\rho<\sqrt5/2$.
\end{proof}

\paragraph{What this does and does not identify.}
The simultaneous construction proves a lower bound on the
four-variable projection dimension, without assuming that
the row--column torus accounts for every generic fiber.
The generic maximal-rank formula
$\min\{\binom{m+3}{3},4m^2-2m+2\}$ is not used or proved.
The lower bound here is enough to retain four-row
multiplicity gaps under an unbounded range of genuine
padding, while the low-row exclusion of G-09 applies at
larger padding. No literature-priority claim is made for
these dimension deductions.

